\section{Ethics and Broader Impacts}

Description-conditioned speech can improve accessibility, multilingual
interfaces, educational tools, and local control over speech infrastructure.
It can also be used to create misleading or harassing audio. Removing
reference-audio cloning from this runtime narrows one misuse path but does not
make generated speech inherently benign: a textual description may still seek
to evoke a recognizable person, role, or demographic stereotype.

Deployers should require authorization, rate limits, abuse monitoring,
appropriate disclosure that audio is synthetic, and a retention policy for
input text and generated audio. Applications involving an identifiable person
require consent and applicable legal review. Provenance or watermarking should
be evaluated where the deployment context supports it; this runtime does not
claim to add an inaudible watermark.

Benchmark publication also has an integrity obligation. Failed trials,
contamination, unsupported languages, unknown EOS state, unavailable sensors,
and negative results must remain visible. Energy measurements from one device
should not be generalized into a lifecycle environmental claim without a
broader accounting boundary.
